
% ----------------------------------------------------------------------
%                   LATEX TEMPLATE FOR PhD THESIS
% ----------------------------------------------------------------------

% based on Harish Bhanderi's PhD/MPhil template, then Uni Cambridge
% http://www-h.eng.cam.ac.uk/help/tpl/textprocessing/ThesisStyle/
% corrected and extended in 2007 by Jakob Suckale, then MPI-CBG PhD programme
% and made available through OpenWetWare.org - the free biology wiki


%: Style file for Latex
% Most style definitions are in the external file PhDthesisPSnPDF.
% In this template package, it can be found in ./Latex/Classes/
\documentclass[twoside,11pt]{Latex/Classes/PhDthesisPSnPDF}


%: Macro file for Latex
% Macros help you summarise frequently repeated Latex commands.
% Here, they are placed in an external file /Latex/Macros/MacroFile1.tex
% An macro that you may use frequently is the figuremacro (see introduction.tex)
% This file contains macros that can be called up from connected TeX files
% It helps to summarise repeated code, e.g. figure insertion (see below).

% insert a centered figure with caption and description
% parameters 1:filename, 2:title, 3:description and label
\newcommand{\figuremacro}[3]{
	\begin{figure}[htbp]
		\centering
		\includegraphics[width=1\textwidth]{#1}
		\caption[#2]{\textbf{#2} - #3}
		\label{#1}
	\end{figure}
}

% insert a centered figure with caption and description AND WIDTH
% parameters 1:filename, 2:title, 3:description and label, 4: textwidth
% textwidth 1 means as text, 0.5 means half the width of the text
\newcommand{\figuremacroW}[4]{
	\begin{figure}[htbp]
		\centering
		\includegraphics[width=#4\textwidth]{#1}
		\caption[#2]{\textbf{#2} - #3}
		\label{#1}
	\end{figure}
}

% inserts a figure with wrapped around text; only suitable for NARROW figs
% o is for outside on a double paged document; others: l, r, i(inside)
% text and figure will each be half of the document width
% note: long captions often crash with adjacent content; take care
% in general: above 2 macro produce more reliable layout
\newcommand{\figuremacroN}[3]{
	\begin{wrapfigure}{o}{0.5\textwidth}
		\centering
		\includegraphics[width=0.48\textwidth]{#1}
		\caption[#2]{{\small\textbf{#2} - #3}}
		\label{#1}
	\end{wrapfigure}
}

% predefined commands by Harish
\newcommand{\PdfPsText}[2]{
  \ifpdf
     #1
  \else
     #2
  \fi
}

\newcommand{\IncludeGraphicsH}[3]{
  \PdfPsText{\includegraphics[height=#2]{#1}}{\includegraphics[bb = #3, height=#2]{#1}}
}

\newcommand{\IncludeGraphicsW}[3]{
  \PdfPsText{\includegraphics[width=#2]{#1}}{\includegraphics[bb = #3, width=#2]{#1}}
}

\newcommand{\InsertFig}[3]{
  \begin{figure}[!htbp]
    \begin{center}
      \leavevmode
      #1
      \caption{#2}
      \label{#3}
    \end{center}
  \end{figure}
}


%%% Local Variables: 
%%% mode: latex
%%% TeX-master: "~/Documents/LaTeX/CUEDThesisPSnPDF/thesis"
%%% End: 




%: ----------------------------------------------------------------------
%:                  TITLE PAGE: name, degree,..
% ----------------------------------------------------------------------
% below is to generate the title page with crest and author name

%if output to PDF then put the following in PDF header
\ifpdf  
    \pdfinfo { /Title  (PhD and MPhil Thesis Classes)
               /Creator (TeX)
               /Producer (pdfTeX)
               /Author (YourName your@email.net)
               /CreationDate (D:YYYYMMDDhhmmss)  %format D:YYYYMMDDhhmmss
               /ModDate (D:YYYYMMDDhhmm)
               /Subject (xyz)
               /Keywords (add, your, keywords, here) }
    \pdfcatalog { /PageMode (/UseOutlines)
                  /OpenAction (fitbh)  }
\fi


\title{Instrumentation}



% ----------------------------------------------------------------------
% The section below defines www links/email for author and institutions
% They will appear on the title page of the PDF and can be clicked
\ifpdf
  \author{\href{mailto:20261142@student.uwa.edu.au}{Beau Trepp}}
%  \cityofbirth{born in XYZ} % uncomment this if your university requires this
%  % If city of birth is required, also uncomment 2 sections in PhDthesisPSnPDF
%  % Just search for the "city" and you'll find them.
  \collegeordept{\href{http://ee.uwa.edu.au}{Electrical, Electronic and Computer Engineering}}
  \university{\href{http://www.uwa.edu.au}{University Of Western Australia}}

  % The crest is a graphics file of the logo of your research institution.
  % Place it in ./0_frontmatter/figures and specify the width
  \crest{\includegraphics[width=4cm]{logo}}
  
% If you are not creating a PDF then use the following. The default is PDF.
\else
  \author{YourName}
%  \cityofbirth{born in XYZ}
  \collegeordept{CollegeOrDept}
  \university{University}
  \crest{\includegraphics[width=4cm]{logo}}
\fi

%\renewcommand{\submittedtext}{change the default text here if needed}
\degree{Bachelor of Computer Science/ Bachelor of Electronic Engineering}
\degreedate{2011 November}


% ----------------------------------------------------------------------
       
% turn of those nasty overfull and underfull hboxes
\hbadness=10000
\hfuzz=50pt


%: --------------------------------------------------------------
%:                  FRONT MATTER: dedications, abstract,..
% --------------------------------------------------------------

\begin{document}

%\language{english}

% sets line spacing
\renewcommand\baselinestretch{1.2}
\baselineskip=18pt plus1pt


%: ----------------------- generate cover page ------------------------

\maketitle  % command to print the title page with above variables


%: ----------------------- cover page back side ------------------------
% Your research institution may require reviewer names, etc.
% This cover back side is required by Dresden Med Fac; uncomment if needed.

\newpage
\vspace{10mm}
1. Reviewer: Name

\vspace{10mm}
2. Reviewer: 

\vspace{20mm}
Day of the defense:

\vspace{20mm}
\hspace{70mm}Signature from head of PhD committee:



%: ----------------------- abstract ------------------------

% Your institution may have specific regulations if you need an abstract and where it is to be placed in the document. The default here is just after title.


% Thesis Abstract -----------------------------------------------------


%\begin{abstractslong}    %uncommenting this line, gives a different abstract heading
\begin{abstracts}        %this creates the heading for the abstract page

The topic of electric vehicles is becoming increasingly popular due to rising fuel costs and growing concern over emissions. Despite this attention, most electric vehicles have little or no telemetry systems, making many aspects of there operation and efficiency a mystery. 

The aim of this project was to develop an extend-able system in order to capture various data-points that can be available in a vehicle, as well as an interface to display this data inside the car. The design developed differs from traditional embedded systems by being completely modular. It uses existing network protocols to allow the system to be distributed between various smaller embedded components. This will enable it to be easily extended, should the need for more data-points arise, and allows the use of many smaller systems to be implemented incrementally, rather than one expensive monolithic design.

In order to help facilitate robustness and code re-use, a windowing toolkit was also developed. This provides a common platform for user interaction, and defines pre-built components in order to speed up the design and implementations of the user interface.

By exposing and recording more data, deeper analysis can be done on the efficiency of the car, and help justify different technological improvements to the vehicle. The higher granularity of data acquired can also be used to analyze the economy of the vehicle in different conditions.

\end{abstracts}
%\end{abstractlongs}


% ---------------------------------------------------------------------- 


% The original template provides and abstractseparate environment, if your institution requires them to be separate. I think it's easier to print the abstract from the complete thesis by restricting printing to the relevant page.
% \begin{abstractseparate}
%   
% Thesis Abstract -----------------------------------------------------


%\begin{abstractslong}    %uncommenting this line, gives a different abstract heading
\begin{abstracts}        %this creates the heading for the abstract page

The topic of electric vehicles is becoming increasingly popular due to rising fuel costs and growing concern over emissions. Despite this attention, most electric vehicles have little or no telemetry systems, making many aspects of there operation and efficiency a mystery. 

The aim of this project was to develop an extend-able system in order to capture various data-points that can be available in a vehicle, as well as an interface to display this data inside the car. The design developed differs from traditional embedded systems by being completely modular. It uses existing network protocols to allow the system to be distributed between various smaller embedded components. This will enable it to be easily extended, should the need for more data-points arise, and allows the use of many smaller systems to be implemented incrementally, rather than one expensive monolithic design.

In order to help facilitate robustness and code re-use, a windowing toolkit was also developed. This provides a common platform for user interaction, and defines pre-built components in order to speed up the design and implementations of the user interface.

By exposing and recording more data, deeper analysis can be done on the efficiency of the car, and help justify different technological improvements to the vehicle. The higher granularity of data acquired can also be used to analyze the economy of the vehicle in different conditions.

\end{abstracts}
%\end{abstractlongs}


% ---------------------------------------------------------------------- 

% \end{abstractseparate}


%: ----------------------- tie in front matter ------------------------

\frontmatter
% Thesis Dedictation ---------------------------------------------------

\begin{dedication} %this creates the heading for the dedication page

%To ...

\begin{flushright}
 

Beau Trepp 

28 Waterfall Rd

Forrestfield WA 6058 
\end{flushright}

\begin{flushleft}
The Dean
Faculty of Engineering Computing and Mathematics \newline
The University of Western Australia \newline
35 Stirling Highway \newline
CRAWLEY WA 6009 
\end{flushleft}

\begin{flushleft}
Dear Sir, \newline
  \newline
I submit to you this dissertation entitled ― "Design of an embedded data acquisition and display system using a modular approach" in partial fulfillment of the requirement of the award of Bachelor of Engineering/Bachelor of Computer Science.
 \newline \newline \newline
Yours Faithfully, \newline
\newline \newline
Beau Trepp \newline
20261142
\end{flushleft}
\end{dedication}

% ----------------------------------------------------------------------
% Thesis Acknowledgements ------------------------------------------------


%\begin{acknowledgementslong} %uncommenting this line, gives a different acknowledgements heading
\begin{acknowledgements}      %this creates the heading for the acknowlegments

During work on this project, I received help from many people whom I would like to thank.

The members of the gumstix and ZeroMQ mailing lists who helped me with issues regarding compilation of both code-bases.

Ian Hooper who helped with work on the vehicle, and provided knowledge into the construction of the vehicle.

Jonathan Oakley who showed helped me with understanding the safety issues regarding driving the car, and for showing me the routes used to performance test it.

The maintainers of the open icon library, for the images used in the project.

To my supervisor Prof Thomas Bräunl for the opportunity to work on the project and vehicle and for guidance during the project.

Kevin Vinsen, for looking after the REV project while Thomas was on sabbatical.

Finally I would like to thank all the members of the REV Team, for providing assistance and support throughout the project.

\end{acknowledgements}
%\end{acknowledgmentslong}

% ------------------------------------------------------------------------





%: ----------------------- contents ------------------------

\setcounter{secnumdepth}{3} % organisational level that receives a numbers
\setcounter{tocdepth}{3}    % print table of contents for level 3
\tableofcontents            % print the table of contents
% levels are: 0 - chapter, 1 - section, 2 - subsection, 3 - subsection


%: ----------------------- list of figures/tables ------------------------

\listoffigures	% print list of figures

\listoftables  % print list of tables


%: ----------------------- glossary ------------------------

% Tie in external source file for definitions: /0_frontmatter/glossary.tex
% Glossary entries can also be defined in the main text. See glossary.tex
% this file is called up by thesis.tex
% content in this file will be fed into the main document

% Glossary entries are defined with the command \nomenclature{1}{2}
% 1 = Entry name, e.g. abbreviation; 2 = Explanation
% You can place all explanations in this separate file or declare them in the middle of the text. Either way they will be collected in the glossary.

% required to print nomenclature name to page header
\markboth{\MakeUppercase{\nomname}}{\MakeUppercase{\nomname}}


% ----------------------- contents from here ------------------------

% chemicals
\nomenclature{DAPI}{4',6-diamidino-2-phenylindole; a fluorescent stain that binds strongly to DNA and serves to marks the nucleus in fluorescence microscopy} 
\nomenclature{DEPC}{diethyl-pyro-carbonate; used to remove RNA-degrading enzymes (RNAases) from water and laboratory utensils}
\nomenclature{DMSO}{dimethyl sulfoxide; organic solvent, readily passes through skin, cryoprotectant in cell culture}
\nomenclature{EDTA}{Ethylene-diamine-tetraacetic acid; a chelating (two-pronged) molecule used to sequester most divalent (or trivalent) metal ions, such as calcium (Ca$^{2+}$) and magnesium (Mg$^{2+}$), copper (Cu$^{2+}$), or iron (Fe$^{2+}$ / Fe$^{3+}$)}



 

\begin{multicols}{2} % \begin{multicols}{#columns}[header text][space]
\begin{footnotesize} % scriptsize(7) < footnotesize(8) < small (9) < normal (10)

\printnomenclature[1.5cm] % [] = distance between entry and description
\label{nom} % target name for links to glossary

\end{footnotesize}
\end{multicols}



%: --------------------------------------------------------------
%:                  MAIN DOCUMENT SECTION
% --------------------------------------------------------------

% the main text starts here with the introduction, 1st chapter,...
\mainmatter

\renewcommand{\chaptername}{} % uncomment to print only "1" not "Chapter 1"


%: ----------------------- subdocuments ------------------------

% Parts of the thesis are included below. Rename the files as required.
% But take care that the paths match. You can also change the order of appearance by moving the include commands.


% this file is called up by thesis.tex
% content in this file will be fed into the main document

%: ----------------------- introduction file header -----------------------
\chapter{Introduction}

% the code below specifies where the figures are stored
\ifpdf
    \graphicspath{{1_introduction/figures/PNG/}{1_introduction/figures/PDF/}{1_introduction/figures/}}
\else
    \graphicspath{{1_introduction/figures/EPS/}{1_introduction/figures/}}
\fi

% ----------------------------------------------------------------------
%: ----------------------- introduction content ----------------------- 
% ----------------------------------------------------------------------



%: ----------------------- HELP: latex document organisation
% the commands below help you to subdivide and organise your thesis
%    \chapter{}       = level 1, top level
%    \section{}       = level 2
%    \subsection{}    = level 3
%    \subsubsection{} = level 4
% note that everything after the percentage sign is hidden from output



\section{put section name here} % section headings are printed smaller than chapter names
% intro
Write your text without any further commands, like this:.... Any organised system requires energy, be it a machine of some kind or a live organism. Energy is needed to win the uphill battle against entropy and pull together lifeless molecules to be able to do something in this world, like complete a PhD.



\subsection{Name your subsection} % subsection headings are again smaller than section names
% lead
Different organised systems have different energy currencies. The machines that enable us to do science like sizzling electricity but at a controlled voltage. Earth's living beings are no different, except that they have developed another preference. They thrive on various chemicals. 

% dextran, starch, glycogen
Most organisms use polymers of glucose units for energy storage and differ only slightly in the way they link together monomers to sometimes gigantic macromolecules. Dextran of bacteria is made from long chains of $\alpha$-1,6-linked glucose units. 

%: ----------------------- HELP: special characters
% above you can see how special characters are coded; e.g. $\alpha$
% below are the most frequently used codes:
%$\alpha$  $\beta$  $\gamma$  $\delta$

%$^{chars to be superscripted}$  OR $^x$ (for a single character)
%$_{chars to be suberscripted}$  OR $_x$

%>  $>$  greater,  <  $<$  less
%≥  $\ge$  greater than or equal, ≤  $\ge$  lesser than or equal
%~  $\sim$  similar to

%$^{\circ}$C   ° as in degree C
%±  \pm     plus/minus sign

%$\AA$     produces  Å (Angstrom)




% dextran, starch, glycogen continued
Starch of plants and glycogen of animals consists of $\alpha$-1,4-glycosidic glucose polymers \cite{lastname07}. See figure \ref{largepotato} for a comparison of glucose polymer structure and chemistry. 

Two references can be placed separated by a comma \cite{lastname07,name06}.

%: ----------------------- HELP: references
% References can be links to figures, tables, sections, or references.
% For figures, tables, and text you define the target of the link with \label{XYZ}. Then you call cross-link with the command \ref{XYZ}, as above
% Citations are bound in a very similar way with \cite{XYZ}. You store your references in a BibTex file with a programme like BibDesk.





\figuremacro{largepotato}{A common glucose polymers}{The figure shows starch granules in potato cells, taken from \href{http://molecularexpressions.com/micro/gallery/burgersnfries/burgersnfries4.html}{Molecular Expressions}.}

%: ----------------------- HELP: adding figures with macros
% This template provides a very convenient way to add figures with minimal code.
% \figuremacro{1}{2}{3}{4} calls up a series of commands formating your image.
% 1 = name of the file without extension; PNG, JPEG is ok; GIF doesn't work
% 2 = title of the figure AND the name of the label for cross-linking
% 3 = caption text for the figure

%: ----------------------- HELP: www links
% You can also see above how, www links are placed
% \href{http://www.something.net}{link text}

\figuremacroW{largepotato}{Title}{Caption}{0.8}
% variation of the above macro with a width setting
% \figuremacroW{1}{2}{3}{4}
% 1-3 as above
% 4 = size relative to text width which is 1; use this to reduce figures




Insulin stimulates the following processes:

\begin{itemize}
\item muscle and fat cells remove glucose from the blood,
\item cells breakdown glucose via glycolysis and the citrate cycle, storing its energy in the form of ATP,
\item liver and muscle store glucose as glycogen as a short-term energy reserve,
\item adipose tissue stores glucose as fat for long-term energy reserve, and
\item cells use glucose for protein synthesis.
\end{itemize}

%: ----------------------- HELP: lists
% This is how you generate lists in LaTeX.
% If you replace {itemize} by {enumerate} you get a numbered list.


 


%: ----------------------- HELP: tables
% Directly coding tables in latex is tiresome. See below.
% I would recommend using a converter macro that allows you to make the table in Excel and convert them into latex code which you can then paste into your doc.
% This is the link: http://www.softpedia.com/get/Office-tools/Other-Office-Tools/Excel2Latex.shtml
% It's a Excel template file containing a macro for the conversion.

\begin{table}[htdp]
\centering
\begin{tabular}{ccc} % ccc means 3 columns, all centered; alternatives are l, r

{\bf Gene} & {\bf GeneID} & {\bf Length} \\ 
% & denotes the end of a cell/column, \\ changes to next table row
\hline % draws a line under the column headers

human latexin & 1234 & 14.9 kbps \\
mouse latexin & 2345 & 10.1 kbps \\
rat latexin   & 3456 & 9.6 kbps \\
% Watch out. Every line must have 3 columns = 2x &. 
% Otherwise you will get an error.

\end{tabular}
\caption[title of table]{\textbf{title of table} - Overview of latexin genes.}
% You only need to write the title twice if you don't want it to appear in bold in the list of tables.
\label{latexin_genes} % label for cross-links with \ref{latexin_genes}
\end{table}



% There you go. You already know the most important things.


% ----------------------------------------------------------------------



	% background information
% this file is called up by thesis.tex
% content in this file will be fed into the main document

\chapter{Aims of the project} % top level followed by section, subsection


% ----------------------- contents from here ------------------------

\section{Final Aim}

The Ultimate goal of the project is to developed a distributed system for use in an automotive environment. This will culminate into a completed system, provided data logging functionality and a user interface to view the live data.

\section{Goals}

The goals this project aimed to achieve were set out as follows.

\begin{enumerate}
\item  Investigate the use messaging protocols inside a minimal embedded systems
\item  Develop GPS capability
\item  Develop Battery monitor capabilty
\item  Integrate this data into a user display
\item  Log this data to be reviewed later
\end{enumerate}

During development, the author also kept the following goals in-line, to make the system easier to develop for.  This was done to give the software a better chance of being developed further by future projects on this topic.


\begin{enumerate}
\item  Allow the system to be easily extend-able
\item  Attempt to minimize future exposure to concurrency
\item  Use simpler individual designs, to make it easier for engineers (as opposed to programmers) to modify
\item Develop standards for inter-module communication
\item Develop common UI elements to facilitate ease-of-layout and interface design.
\item Develop the system in such a way that it is portable to other devices. The hardware used is slightly older, and the system will be developed in a way that it can be ported across to new hardware.
\end{enumerate}


% ---------------------------------------------------------------------------
% ----------------------- end of thesis sub-document ------------------------
% ---------------------------------------------------------------------------						% aims of the project

\include{3/XYZ}			
\include{4/XYZ}	
\include{5/XYZ}
\include{6/XYZ}

% this file is called up by thesis.tex
% content in this file will be fed into the main document

\chapter{Performance} % top level followed by section, subsection


% ----------------------- paths to graphics ------------------------

% change according to folder and file names
\ifpdf
    \graphicspath{{7/figures/PNG/}{7/figures/PDF/}{7/figures/}}
\else
    \graphicspath{{7/figures/EPS/}{7/figures/}}
\fi


% ----------------------- contents from here ------------------------


\section{Limitations}

In testing the final system, issues arose which limited the depth in which the systems could be evaluated. As such only partial evaluation has been undertaken. This was done in an effort to provide some indication as to the stability of the system, and this section will provide analysis over what data could be obtained.

\subsection{Restricted Access to Vehicle}
\label{sec:access}
On the 15th of August the vehicle that was the main target for development was involved in a car accident. For safety reasons, the author was not allowed to work on the car in any form after this accident. This severely limited the amount of testing that could take place, as many aspects of the hardware are installed in the car which was no longer accessible.

\section{Message Timings}

This section reveals the time in between messages for the various components that transmit them. This was done using an eyebot to run the programs, and the messages were received on another laptop. Errors in the measurement can be contributed to the process switching active on both devices, as well as the fact that the laptop was communicating via a wireless network connection. All tests were conducted for approximately 20 minutes.

\subsection{Arduino Board}

Figure \ref{arduinotiming} shows the time between messages received vs the running time in the program. It is interesting to note that there are occasions when the time between two messages can be as much as 4 seconds. As the graph displays, this phenomenon appears to occur every 7 seconds. This is most likely caused by delays in the operating system on the receiving or sending end, as the process will not switch to the foreground automatically. As the delay is rarely more than 4 seconds, this will not interfere with the timeouts implemented in section \ref{toolkitfilter}

\figuremacro{arduinotiming}{Inter message timing for Arduino network messages}{The time delay between messages plotted against the time the program has been running for}

\subsection{Accelerometer}

The accelerometer runs at a much faster rate than the Arduino board. This is reflected in graph \ref{imutiming}. It can easily be seen that this program sends out messages much more frequently.  It still suffers the same problem with longer delays occurring. However each occurrence is still well under any timeout limits. 

\figuremacro{imutiming}{Inter message timing for accelerometer network messages}{The time delay between messages plotted against the time the program has been running for}

\subsection{GPS and Battery Monitor}

Due to both being physically installed in the vehicle, it was not possible to evaluate the message timings of these systems. See section \ref{sec:access}


\section{Memory Utilization}

This section shows plots of the memory usage of each daemon vs the running time of the program. If a program has a memory leak, it will eventually cause exhaust the systems memory and cause instability. These graphs exist to highlight any issues that may occur with running programs.

\subsection{Arduino Board}

Sadly, as shown in figure \ref{memoryarduino}, the Arduino based daemon does have a minor memory leak. After running the program for one hour, the memory usage had increased by 100kb. This memory leak is not too much of a problem. The system will never be on for more than a day at one time. After a day running non-stop, the program will only consume a little over 2mb. This is nothing compared to the 64mb of ram available to the system, so this problem is very minor.


\figuremacro{memoryarduino}{Memory usage of the Arduino daemon}{Memory usage (in kb) plotted against the time the program has been running for}

\subsection{Accelerometer}

As shown in figure \ref{memoryimu} the Accelerometer daemon has a memory leak at the rate of 760kb per hour. This is much worse than the leak in the Arduino board. Extrapolating this usage out to one day yields a memory usage of approximately 18mb. This is still small enough to not cause major problems, though it does warrant an investigation into the cause of the memory leak. As long as the system is not run for over a day without a restart, this leak will not be a problem.

\figuremacroW{memoryimu}{Memory usage of the Accelerometer daemon}{Memory usage (in kb) plotted against the time the program has been running for}{0.8}

\subsection{User interface}

The user interface was tested with a simulator sending messages from all four components developed in section \ref{sec:design}. As figure \ref{memorygui} shows, the user interface does not have any memory leaks. The interface can run for any amount of time and will not cause system instability. This is a good result as the user interface will always remain responsive.

\figuremacroW{memorygui}{Memory usage of the user interface}{Memory usage (in kb) plotted against the time the program has been running for}{0.8}

\subsection{GPS and Battery Monitor}

Due to both being physically installed in the vehicle, it was not possible to evaluate the message timings of these systems. See section \ref{sec:access}. Running these programs without access to the hardware would not yield reliable results.




% ---------------------------------------------------------------------------
% ----------------------- end of thesis sub-document ------------------------
% ---------------------------------------------------------------------------               % discussion of results


% this file is called up by thesis.tex
% content in this file will be fed into the main document

\chapter{Conclusions} % top level followed by section, subsection


% ----------------------- paths to graphics ------------------------

% change according to folder and file names
\ifpdf
    \graphicspath{{8/figures/PNG/}{8/figures/PDF/}{8/figures/}}
\else
    \graphicspath{{8/figures/EPS/}{8/figures/}}
\fi

% ----------------------- contents from here ------------------------



This project shows that it is possible to use high level programming techniques when developing for embedded devices. The hardware used in this project was powerful enough to separate into multiple components. This is significant as it provides flexibility in the development of embedded devices and highlights how these products are becoming more like traditional computers.

This project also illustrates the viability of using existing networking techniques, such as TCP/IP, in providing communication between different devices. The application of this technology to embedded systems provides increased robustness and decreases development time.  This shows that a network of small embedded systems that each contribute a portion of the overall system functionality is a viable solution to the problems of scaling and system cost.

The ability to automatically log all the information recorded by the car is of great help to the goals of the REV team. The group will be able to log much more detailed data, and the system has been constructed in a way to make it easily extend-able. This will ultimately help in the research and development of the vehicle the system is running on.

\section{Limitations}

The system currently has a few limitations. These are caused by various constraints in the hardware and the design of the system.

The system will be unable to support too many more additions with it's current processing power. The CPU is not able to run many high frequency data sources at the same time. This will limit the future expansion unless new hardware is acquired.

The system also suffers some issues with the memory utilization of some components. Further investigation will need to take place in order to eliminate the bugs that are causing these issues.

Due to time constraints the Arduino and Accelerometer were not developed fully in regards to the user interface. This is a limitation that can be overcome in the future, when a more specific use of these components is required.

\section{Advantages}

The system is easily expendable, providing a framework for the transmission and display of any data that the group may want to develop in the future. The development of the windowing toolkit makes it much easier to work with the touch interface on the device. 

The messaging protocol developed, and the use of a message abstraction, makes it much easier to communicate between different sections. The fact that the complicated components of networking are abstracted from the developer make it easier to conceptualize what is happening. A component in this system only needs to worry about the data being sent, and not a complicated method of ensuring it is received correctly.

Due to it's distributed nature, the system is easy and cost-effective to extend. Rather than having to source expensive new hardware to perform all the tasks, cheaper hardware can be used to take the role of some of the tasks. This allows the complete system to adapt to the conditions that it is being used in.

\section{Applications}

The system has applications in the project it was developed for. It provides a clean user interface for the display of data, and logging functionality so it can be analyzed later.

The system is a proof of concept that embedded systems can communicate reliably using network protocols. This concept could be applied to newly developed systems, enabling products that can be purchased in sections, with each section providing different or improved functionality. This paves the way for embedded systems to be configured via pre-built components , rather than developed using one monolithic structure. 

\section{Future work}

The first aspect of future work is the finishing of the Arduino and Accelerometer components. These components are able to be decoded, but do not have very useful user interfaces. 

Currently the networking is not encrypted, which does expose the possibility that it can be attacked. This does require access to the network, which is not easily, and understanding of the protocols being communicated. Future work would include the investigation of SSL in order to prevent any third party from viewing the data that is transmitted.

The windowing toolkit could be further developed by providing rendering of the full alphabet at varying sizes. This would simplify the use of the toolkit, as it would be much easier to make text appear. The toolkit could also be further developed into it's own standalone package, either with dependencies on the existing libraries or not.

Navigation could be added to the maps portion. A path-finding algorithm could be developed and implemented that would be able to direct the driver towards a specific destination.

Further research could be conducted into sensor fusion. As the system has access to GPS, the speedometer pulse frequency, current, as well as the acceleration the car is undertaking a system could be developed in order to compensate for loss of a sensor. For instance the other sensors could approximate the GPS position when inside a tunnel. This would increase the user experience, has it would appear that the GPS works in all locations.

New hardware should also be investigated for the system. Either in the form of additional devices to offset the load, or replacement of the main component running the system. This would increase the processing power of the system, and allow more complicated functionality to be developed.


% ---------------------------------------------------------------------------
%: ----------------------- end of thesis sub-document ------------------------
% ---------------------------------------------------------------------------



 






        % description of lab methods




% --------------------------------------------------------------
%:                  BACK MATTER: appendices, refs,..
% --------------------------------------------------------------

% the back matter: appendix and references close the thesis


%: ----------------------- bibliography ------------------------

% The section below defines how references are listed and formatted
% The default below is 2 columns, small font, complete author names.
% Entries are also linked back to the page number in the text and to external URL if provided in the BibTex file.

% PhDbiblio-url2 = names small caps, title bold & hyperlinked, link to page 
\begin{multicols}{2} % \begin{multicols}{ # columns}[ header text][ space]
\begin{tiny} % tiny(5) < scriptsize(7) < footnotesize(8) < small (9)

\bibliographystyle{Latex/Classes/PhDbiblio-url2} % Title is link if provided
\renewcommand{\bibname}{References} % changes the header; default: Bibliography

\bibliography{9_backmatter/references} % adjust this to fit your BibTex file

\end{tiny}
\end{multicols}

% --------------------------------------------------------------
% Various bibliography styles exit. Replace above style as desired.

% in-text refs: (1) (1; 2)
% ref list: alphabetical; author(s) in small caps; initials last name; page(s)
%\bibliographystyle{Latex/Classes/PhDbiblio-case} % title forced lower case
%\bibliographystyle{Latex/Classes/PhDbiblio-bold} % title as in bibtex but bold
%\bibliographystyle{Latex/Classes/PhDbiblio-url} % bold + www link if provided

%\bibliographystyle{Latex/Classes/jmb} % calls style file jmb.bst
% in-text refs: author (year) without brackets
% ref list: alphabetical; author(s) in normal font; last name, initials; page(s)

%\bibliographystyle{plainnat} % calls style file plainnat.bst
% in-text refs: author (year) without brackets
% (this works with package natbib)


% --------------------------------------------------------------

% according to Dresden med fac summary has to be at the end
%
% Thesis Abstract -----------------------------------------------------


%\begin{abstractslong}    %uncommenting this line, gives a different abstract heading
\begin{abstracts}        %this creates the heading for the abstract page

The topic of electric vehicles is becoming increasingly popular due to rising fuel costs and growing concern over emissions. Despite this attention, most electric vehicles have little or no telemetry systems, making many aspects of there operation and efficiency a mystery. 

The aim of this project was to develop an extend-able system in order to capture various data-points that can be available in a vehicle, as well as an interface to display this data inside the car. The design developed differs from traditional embedded systems by being completely modular. It uses existing network protocols to allow the system to be distributed between various smaller embedded components. This will enable it to be easily extended, should the need for more data-points arise, and allows the use of many smaller systems to be implemented incrementally, rather than one expensive monolithic design.

In order to help facilitate robustness and code re-use, a windowing toolkit was also developed. This provides a common platform for user interaction, and defines pre-built components in order to speed up the design and implementations of the user interface.

By exposing and recording more data, deeper analysis can be done on the efficiency of the car, and help justify different technological improvements to the vehicle. The higher granularity of data acquired can also be used to analyze the economy of the vehicle in different conditions.

\end{abstracts}
%\end{abstractlongs}


% ---------------------------------------------------------------------- 


%: Declaration of originality

% Thesis statement of originality -------------------------------------

% Depending on the regulations of your faculty you may need a declaration like the one below. This specific one is from the medical faculty of the university of Dresden.

\begin{declaration}        %this creates the heading for the declaration page

I herewith declare that I have produced this paper without the prohibited assistance of third parties and without making use of aids other than those specified; notions taken over directly or indirectly from other sources have been identified as such. This paper has not previously been presented in identical or similar form to any other German or foreign examination board.

The thesis work was conducted from XXX to YYY under the supervision of PI at ZZZ.

\vspace{10mm}

CITY,


\end{declaration}


% ----------------------------------------------------------------------



\end{document}
